\section{Wind versus current: what GPS and a compass can separate}
\label{sec:cond-observe}

The pure windage estimator fits heading-dependent neutral drift. It determines
a constant vector and a double-angle term, without separating wind and current.

\subsection{A model of a drifting boat}

The following leeway model motivates the harmonic fit; the estimator does not
solve an air--water drag balance. Let $\bm c$ be the current over the ground and $\bm w$ the wind \emph{over the
water}, the air's velocity relative to the surface the hull sits in. With dead
thrusters the steady boat velocity over the water is set by a balance between air
drag on the topsides and water drag on the hull.

\paragraph{One axis first.} Along a single body axis, with both drags quadratic,
$c_a\,|w-u|(w-u) = c_h\,|u|\,u$ where $u$ is the boat's speed over the water and $c_a$,
$c_h$ (air, hull) lump density, area and drag coefficient for air and water. With $u$ between 0
and $w$ both sides have the same sign, so taking square roots,
\[
  \sqrt{c_a}\,(w-u) = \sqrt{c_h}\,u \quad\Longrightarrow\quad
  u = \frac{\sqrt{c_a/c_h}}{1 + \sqrt{c_a/c_h}}\;w \equiv k\,w .
\]
The leeway is \emph{linear} in the wind, with a coefficient $k$ that depends only
on the ratio of areas and coefficients (a few per cent for a small boat). That
linearity needs both drags to have the same power law, which is why it holds for
quadratic air drag on quadratic water drag. (Had we written the wind over the
ground, the leeway would be $k(\bm w_\text{ground} - \bm c)$: the difference is $k\bm c$,
about $0.003$\,m/s for $k = 0.03$ and a $0.1$\,m/s current, and it does not change
anything below.)

\paragraph{Two axes.} A catamaran presents far more area abeam than ahead, so the
leeway coefficient differs by body axis. Linearized, with a symmetric hull, the
body-frame leeway matrix is $K_b = \operatorname{diag}(k_x, k_y)$. To use it, rotate
the wind into the body, apply $K_b$, and rotate back:
\[
  \bm d(\psi) = \bm c + R(\psi)\,K_b\,R(\psi)^\top\,\bm w .
\]

\paragraph{The key identity.} Write $\bar k = (k_x+k_y)/2$ and
$\Delta k = (k_x - k_y)/2$, so $K_b = \bar k I + \Delta k\,\operatorname{diag}(1,-1)$.
The isotropic part commutes with the rotation, $R\,\bar kI\,R^\top = \bar k I$. The
other part needs a short calculation with $c = \cos\psi$, $s = \sin\psi$:
\[
  R\begin{bmatrix}1&0\\0&-1\end{bmatrix}R^\top
  = \begin{bmatrix}c&-s\\s&c\end{bmatrix}\begin{bmatrix}c&s\\s&-c\end{bmatrix}
  = \begin{bmatrix}c^2-s^2 & 2cs\\ 2cs & s^2-c^2\end{bmatrix}
  = \begin{bmatrix}\cos2\psi & \sin2\psi\\ \sin2\psi & -\cos2\psi\end{bmatrix}
  \equiv \Qm{2\psi}.
\]
$\Qm{2\psi}$ is a reflection across the line at angle $\psi$; it depends on the
\emph{double} angle because a reflection axis at $\psi$ and at $\psi+180^\circ$ is
the same line. So
\[
  \boxed{\ \bm d(\psi) = \underbrace{(\bm c + \bar k\,\bm w)}_{\bm a}
                        + \Qm{2\psi}\,\underbrace{\Delta k\,\bm w}_{\bm q}\ }
\]
As the boat swings round, the drift traces a circle of radius $|\bm q|$ about
$\bm a$, \emph{twice} per revolution (Figure~\ref{fig:cond-windage}). A bow-on
and a stern-on boat drift alike; a beam-on boat drifts more downwind.

\begin{figure}[ht]
\centering
\begin{tikzpicture}
\begin{axis}[cond axis, width=0.45\linewidth, height=0.45\linewidth, axis equal image, scaled ticks=false, tick label style={/pgf/number format/fixed},
  xlabel={$d_E$ (m/s)}, ylabel={$d_N$ (m/s)},
  xmin=-0.005, xmax=0.1, ymin=-0.005, ymax=0.11,
  legend style={font=\small, draw=none, fill=none}, legend pos=south east,
  grid=major, grid style={gray!20}]
  \addplot[sA, thick] table[x=dE, y=dN] {figures/conditions/windage.dat};
  \addlegendentry{$\bm d(\psi)$, $\psi$ one turn}
  \addplot[only marks, mark=*, mark size=2pt, sB] table[x=dE, y=dN] {figures/conditions/windage_pts.dat};
  \addlegendentry{six survey windows}
  \draw[->, thick, black] (axis cs:0,0) -- (axis cs:0.06,0.08);
  \fill (axis cs:0.06,0.08) circle (1.5pt) node[left, font=\small] {$\bm a$};
  \draw[->, thick, red] (axis cs:0.06,0.08) -- (axis cs:0.0699,0.0899) node[above right, font=\small] {$\bm q$};
\end{axis}
\end{tikzpicture}
\caption{Drift as a function of heading for $\bm a = (0.06, 0.08)$\,m/s and a
3.5\,m/s wind toward the south-west with $\Delta k = -0.004$. The circle has
radius $|\bm q| = 0.014$\,m/s and is traversed twice per turn. Six windows at
headings $60^\circ$ apart sample it at only three distinct values of $2\psi$,
because headings $180^\circ$ apart coincide. Data:
\code{figures/conditions/windage*.dat}.}
\label{fig:cond-windage}
\end{figure}

\subsection{What is observable}

The measurements are window means $\bm m_w \approx \bm a + \Qm{2\psi_w}\bm q$ at known
headings. The unknowns of the physics are six numbers, $\bm c$, $\bm w$ and the
two leeway coefficients; the data determine four.
At one heading only the total $\bm d(\psi)$ is measured; $\bm a$ and
$\bm q$ cannot be separated. With sufficient heading spread the weighted
four-parameter linear fit uses rows $[I_2\ \Qm{2\psi_w}]$ and weights $T_w/\Szero$.
Wind and current remain inseparable without leeway coefficients.

\paragraph{The fit's gates.} At least three usable windows and a spectral level
are required. The code checks rank three of $[1,\cos2\psi_c,\sin2\psi_c]$,
then a smallest singular value of at least 0.1 for the centred sine/cosine
columns divided by $\sqrt W$. A nonfinite or nonpositive fitted standard error
abstains. The covariance uses inverse weighted information, overdispersion
and residual-correlation inflation. The fit is heading-dependent when
$|\bm q|/\sigma_q>3$; otherwise it returns the constant pooled drift.
\code{drift\_at} evaluates the harmonic term only for a heading-dependent fit.

\file{mhseals\_testing/analysis/conditions.py (windage, Windage.drift\_at)}
